\chapter{多因子与计量分层答案}

\begin{enumerate}[leftmargin=*]
  \item 先把预测回归 $r_{i,t+1}=a_t+b_tx_{i,t}+u_{i,t+1}$ 与定价回归 $r_{i,t}=\beta_i^T\lambda+\alpha_i+\varepsilon_{i,t}$ 分开；因果对象还需处理和反事实。IC 只评价预测/排序关系，不能把它改名为风险价格。
  \item Pearson 计算原值协方差，Rank IC 计算秩协方差；对单调非线性或离群点分别给出一个小表即可看到结论分离。并列秩、缺失和 winsorize 规则应写入 reference ledger。
  \item 两期 $0.018,0.022$ 的均值为 $0.020$，样本方差 $8\times10^{-6}$，均值方差 $4\times10^{-6}$，标准误 $0.002$；公式来自 $\bar b=T^{-1}\sum_tb_t$ 和 $\widehat{\operatorname{Var}}(\bar b)=s_b^2/T$。中间的 $b_t-\bar b$ 不能省略，否则无法检查自由度。
  \item 投影矩阵 $P_Z=Z(Z^TZ)^{-1}Z^T$，残差 $x^\perp=(I-P_Z)x$。逐步相乘得 $Z^Tx^\perp=Z^Tx-Z^TZ(Z^TZ)^{-1}Z^Tx=0$；若秩亏，逆不存在，应改用 QR/SVD 并报告奇异值，而不是返回任意伪逆结果。
  \item 三期扩展先在独立账本中计算每期回归斜率、IC、价差收益和均值，再让程序读取 fixture 产生 observed；逐列比较数值和哈希。这样可区分输入/时间错误、回归实现错误和收益账本错误。
  \item 日期相等、共线暴露和伪造 expected 分别触发 time-alignment、ill-conditioned-neutralization 和 ledger failure。每个负例只改一个字段，验证退出码、stderr 和失败位置，不能用统一宽容差吞掉三类错误。
  \item 500 项检验先冻结 p 值族、BH 水平和依赖条件，在未参与选择的时间窗核验 IC 衰减、换手、成本、冲击、借券和容量。即使 BH 通过，也只控制特定统计错误比例，不保证成本后收益。
  \item 全球因子在 A 股失效时，先重算公式和估计对象，再查发布日期/复权/成分/缺失，最后查涨跌停、停牌、T+1、融券和成交；每一步保留可复现 fixture。只有把数学、数据和制度三类错误拆开，才能提出修复。
\end{enumerate}
